\ifdefined\gkver\else\def\gkver{student}\fi
\documentclass[\gkver]{gaokaozhenti}
\begin{document}

\chapter{2026年福建}

\begin{enumerate}
\item
%% number: 1
%% paperName: 2026年福建高考第 1 题 4 分
%% typeId: 1
%% type: 单选题
%% chapter: 原子和原子核
%% point: 天然放射性
%% method: 无
%% score: 4
%% degree: 850
%% duplicateId: 0
%% body:
我国科学家在钍基熔盐堆中成功实现了钍—铀核燃料循环，该成果入选“2021 年度中国科学十大进展”。循环涉及的核反应方程包含：① $\ce{^{232}_{90}Th + X -> ^{233}_{90}Th}$ ② $\ce{^{233}_{90}Th -> ^{233}_{91}Pa + Y + \bar{\nu}_{e}}$ ③ $\ce{^{233}_{91}Pa -> ^{233}_{92}U + Y + \bar{\nu}_{e}}$。式中 $\bar{\nu}_{e}$ 是质量可忽略不计的中性粒子，则\xzanswer{D}

\fourchoices[answer=D]
{$\ce{X}$ 是 $\ce{^{0}_{-1}e}$}
{$\ce{Y}$ 是 $\ce{^{1}_{0}n}$}
{方程②属于 $\alpha$ 衰变}
{方程③属于 $\beta$ 衰变}

%% answer: D
%% memo:
\memoanswer{
A．对反应①，设 $\ce{X}$ 的质量数为 $A$、电荷数为 $Z$，由 $232+A=233$，$90+Z=90$，解得 $A=1$，$Z=0$，因此 $\ce{X}$ 是中子 $\ce{^{1}_{0}n}$，故 A 错误；

B．反应②中 $\bar{\nu}_{e}$ 质量可忽略、电荷为 $0$，设 $\ce{Y}$ 的质量数为 $M$、电荷数为 $N$，由 $233=233+M+0$，$90=91+N+0$，解得 $M=0$，$N=-1$，因此 $\ce{Y}$ 是电子 $\ce{^{0}_{-1}e}$，故 B 错误；

C．$\alpha$ 衰变的本质是放出 $\alpha$ 粒子（$\ce{^{4}_{2}He}$），但反应②放出的是电子，不属于 $\alpha$ 衰变，故 C 错误；

D．反应③中母核质量数不变、电荷数增加 $1$，同时放出电子，符合 $\beta$ 衰变的特征（中子转化为质子放出电子），属于 $\beta$ 衰变，故 D 正确。

故选 D。
}

\item
%% number: 2
%% paperName: 2026年福建高考第 2 题 4 分
%% typeId: 1
%% type: 单选题
%% chapter: 机械波
%% point: 波的图象
%% method: 无
%% score: 4
%% degree: 750
%% duplicateId: 0
%% body:
在一均匀介质中，位于 $x=0$ 和 $x=10\Um$ 处的两个波源 $P$ 和 $Q$ 沿 $y$ 轴方向振动，形成了两列相向传播的简谐横波，波速均为 $2\Ums$。$t=0$ 时，两列波分别传播至 $x=4\Um$ 和 $x=6\Um$ 处，波形图如图所示。则\xzanswer{B}
\onepicture[w=5.73cm]{figs/02.png}
\fourchoices[answer=B]
{两列波的周期不同}
{$t=0$ 时，$x=2\Um$ 处质点振动方向沿 $y$ 轴正方向}
{$t=1\Us$ 时，两列波开始相遇}
{从 $t=0$ 到 $t=1\Us$，$x=8\Um$ 处质点沿 $x$ 轴负方向移动了 $2\Um$}
%% answer: B
%% memo:
\memoanswer{
A．从波形图可得两列波的波长均为 $\lambda=4\Um$。已知波速 $v=2\Ums$，根据周期公式 $T=\frac{\lambda}{v}$，代入数据得 $T=\frac{4}{2}\Us=2\Us$。即两列波的周期相同，均为 $2\Us$，故 A 错误；

B．根据同侧法可知 $t=0$ 时，$x=2\Um$ 处质点振动方向沿 $y$ 轴正方向，故 B 正确；

C．$t=0$ 时，两列波分别传播至 $x=4\Um$ 和 $x=6\Um$ 处，两波前端的间距为 $\Delta x=6\Um-4\Um=2\Um$。两波相向传播，相对速度为 $2v$，相遇时间 $t=\frac{\Delta x}{2v}=\frac{2}{2\times2}\Us=0.5\Us$，故 C 错误；

D．简谐横波传播过程中，介质中的质点仅在平衡位置附近沿 $y$ 轴方向振动，不会随波沿 $x$ 轴方向迁移，因此 $x=8\Um$ 处质点不会沿 $x$ 轴移动，故 D 错误。

故选 B。
}

\item
%% number: 3
%% paperName: 2026年福建高考第 3 题 4 分
%% typeId: 1
%% type: 单选题
%% chapter: 磁场
%% point: 磁场的叠加
%% method: 无
%% score: 4
%% degree: 701
%% duplicateId: 0
%% body:
我国航空母舰“福建舰”的电磁弹射系统装配有直线电机。如图，电机水平轨道两侧对称平行固定有多组相同线圈，通过控制线圈电流可产生特定磁场。$P$ 点为过轨道中轴线的竖直面和线圈组 aa$'$ 中轴线的交点，在相邻线圈组 aa$'$ 和 bb$'$ 中通有图示方向、大小分别为 $2I$ 和 $I$ 的电流时，两组线圈在 $P$ 点产生的磁场的方向为\xzanswer{B}
\onepicture[w=6.05cm]{\includesvg{figs/03}}
\fourchoices[answer=B]
{①方向}
{②方向}
{③方向}
{④方向}
%% answer: B
%% memo:
\memoanswer{
由通电螺线管周围磁场分布特点、线圈放置的对称性以及矢量叠加原理可知，线圈组 aa$'$ 在 $P$ 点产生的磁场方向为②方向，而线圈组 bb$'$ 在 $P$ 点产生的磁感应强度方向为④方向。由于 bb$'$ 中电流仅为 aa$'$ 中电流的一半，且其与 $P$ 点的距离更远，由场的叠加原理可知，$P$ 点磁感应强度由 aa$'$ 产生的磁场主导，故两组线圈在 $P$ 点产生的磁场的方向为②方向。

故选 B。
}

\item
%% number: 4
%% paperName: 2026年福建高考第 4 题 4 分
%% typeId: 1
%% type: 单选题
%% chapter: 电场
%% point: 带电粒子的偏转
%% method: 无
%% score: 4
%% degree: 620
%% duplicateId: 0
%% body:
反射飞行时间质谱仪的核心部件中存在真空无场漂移区和匀强电场反射区，竖直剖面如图所示。某次调试过程中，比荷（电荷量与质量之比）均为 $K$ 的离子 a、b 分别以不同速率 $v_a$、$v_b$ 从入口 $O$ 点入射，入射方向与水平方向的夹角均为 $\theta$，经漂移区进入反射区后，在水平向右的电场作用下返回漂移区，最终到达 $O$ 点正上方的探测屏。已知漂移区长度为 $L$，离子运动轨迹处于与电场方向平行的竖直面内，重力不计。要使离子 a、b 从 $O$ 点到探测屏的飞行时间相等，反射区电场强度大小 $E$ 应调为\xzanswer{C}
\onepicture[w=6.84cm]{figs/04.png}
\fourchoices[answer=C]
{$\frac{v_av_b\cos\theta}{KL}$}
{$\frac{2v_av_b\cos\theta}{KL}$}
{$\frac{v_av_b\cos^2\theta}{KL}$}
{$\frac{2v_av_b\cos^2\theta}{KL}$}
%% answer: C
%% memo:
\memoanswer{
设离子 a、b 的质量分别为 $m_a$、$m_b$，电荷量分别为 $q_a$、$q_b$，比荷 $\frac{q_a}{m_a}=\frac{q_b}{m_b}=K$。

a 在漂移区内做匀速直线运动，运动时长 $t_1$ 满足 $v_a\cos\theta\cdot t_1=L$，即 $t_1=\frac{L}{v_a\cos\theta}$。

a 在反射区内做匀减速直线运动，加速度大小 $a=\frac{q_aE}{m_a}=KE$，运动至电场最深处的时长 $t_2$ 满足 $KE\cdot t_2=v_a\cos\theta$，即 $t_2=\frac{v_a\cos\theta}{KE}$。

由运动的对称性可知，a 在质谱仪内的飞行总时间 $t_a=2(t_1+t_2)=2\left(\frac{L}{v_a\cos\theta}+\frac{v_a\cos\theta}{KE}\right)$。

同理，b 的飞行总时间 $t_b=2\left(\frac{L}{v_b\cos\theta}+\frac{v_b\cos\theta}{KE}\right)$。

由题意 $t_a=t_b$，联立可得 $\frac{L}{v_a\cos\theta}+\frac{v_a\cos\theta}{KE}=\frac{L}{v_b\cos\theta}+\frac{v_b\cos\theta}{KE}$，

整理得 $\frac{L(v_b-v_a)}{v_av_b\cos\theta}=\frac{(v_a-v_b)\cos\theta}{KE}$，即 $E=\frac{v_av_b\cos^2\theta}{KL}$。

故选 C。
}

\item
%% number: 5
%% paperName: 2026年福建高考第 5 题 6 分
%% typeId: 2
%% type: 多选题
%% chapter: 运动和力
%% point: 牛顿定律中的图象
%% method: 无
%% score: 6
%% degree: 851
%% duplicateId: 0
%% body:
高速运行的磁悬浮列车利用空气阻力进行制动。在某次测试中，关闭动力后，列车在水平直轨道上受阻力 $f$ 作用做减速运动，$f$ 随时间 $t$ 的变化关系如图所示。则列车在 $0\sim200\Us$ 内\xzanswer{BD}
\onepicture[w=4.93cm]{figs/05.png}
\fourchoices[answer=BD]
{做匀减速直线运动}
{加速度大小随时间不断减小}
{所受阻力 $f$ 的功率随时间不断增大}
{所受阻力 $f$ 的功率随时间不断减小}
%% answer: BD
%% memo:
\memoanswer{
AB．根据牛顿第二定律 $f=ma$ 结合题图可知列车所受阻力逐渐减小，则列车在 $0\sim200\Us$ 内加速度大小随时间不断减小，列车做加速度减小的减速运动，故 A 错误，B 正确；

CD．列车速度减小，摩擦力也不断减小，根据 $P=fv$ 可知列车所受阻力 $f$ 的功率随时间不断减小，故 C 错误，D 正确。

故选 BD。
}

\item
%% number: 6
%% paperName: 2026年福建高考第 6 题 6 分
%% typeId: 2
%% type: 多选题
%% chapter: 曲线运动
%% point: 人造地球卫星
%% method: 无
%% score: 6
%% degree: 670
%% duplicateId: 0
%% body:
1970 年，我国第一颗人造卫星“东方红一号”成功进入绕地椭圆轨道，轨道半长轴约为 $7.8\times10^3\Ukm$。在稀薄大气阻力影响下，卫星运行轨道半长轴缓慢减小，2026 年其值约为 $7.6\times10^3\Ukm$。如图，1970 年和 2026 年的轨道分别记为轨道 Ⅰ 和 Ⅱ，且视作稳定运行轨道。忽略卫星质量变化，则卫星\xzanswer{AC}
\onepicture[w=6.59cm]{figs/06.png}
\fourchoices[answer=AC]
{在轨道 Ⅰ 上的机械能比在轨道 Ⅱ 上的大}
{在轨道 Ⅰ、Ⅱ 上的机械能变化量等于在两轨道远地点的动能变化量}
{在轨道 Ⅰ、Ⅱ 上运行周期之比约为 $39^{\frac{3}{2}}:38^{\frac{3}{2}}$}
{在轨道 Ⅰ、Ⅱ 上远地点的加速度大小之比约为 $39^2:38^2$}
%% answer: AC
%% memo:
\memoanswer{
A．卫星受稀薄大气阻力作用，阻力持续做负功，卫星机械能不断减小；轨道 Ⅰ 的半长轴更大、是更早的轨道，因此卫星在轨道 Ⅰ 上的机械能比在轨道 Ⅱ 上的大，故 A 正确；

B．卫星的机械能是动能与引力势能的总和，两轨道的机械能变化量等于远地点动能变化量与引力势能变化量之和，故 B 错误；

C．根据开普勒第三定律 $\frac{T_1^2}{T_2^2}=\frac{a_1^3}{a_2^3}$，代入 $a_1=7.8\times10^3\Ukm$、$a_2=7.6\times10^3\Ukm$，可得 $\frac{T_1}{T_2}=\left(\frac{7.8}{7.6}\right)^{\frac{3}{2}}=\frac{39^{\frac{3}{2}}}{38^{\frac{3}{2}}}$，故 C 正确；

D．卫星的加速度 $a=\frac{GM}{r^2}$，与到地心距离的平方成反比；轨道 Ⅰ、Ⅱ 的远地点到地心的距离之比不等于 $39:38$，因此加速度之比不等于 $39^2:38^2$，故 D 错误。

故选 AC。
}

\item
%% number: 7
%% paperName: 2026年福建高考第 7 题 6 分
%% typeId: 2
%% type: 多选题
%% chapter: 运动和力
%% point: 动量和能量
%% method: 无
%% score: 6
%% degree: 401
%% duplicateId: 0
%% body:
如图，一倾角 $\theta$ 为 $30^\circ$ 的足够长斜面固定在水平面上，斜面上放置两个由轻质弹簧连接的物块 M 和 N。M、N 的质量分别为 $0.5\Ukg$ 和 $1\Ukg$，M 与斜面的接触面光滑，N 与斜面间的动摩擦因数为 $\frac{\sqrt3}{2}$。初始时，弹簧处于原长状态，M、N 由静止释放。弹簧始终在弹性限度内，最大静摩擦力等于滑动摩擦力。则\xzanswer{AD}
\onepicture[w=4.71cm]{\includesvg{figs/07}}
\fourchoices[answer=AD]
{N 开始运动的瞬间，M 的加速度为 $0$}
{N 的速度达到最大时，M 的速度为 $0$}
{N 开始运动后，M、N 的总动能与弹簧弹性势能之和保持不变}
{M、N 第一次共速时的弹簧弹性势能和第二次共速时的不同}
%% answer: AD
%% memo:
\memoanswer{
A．N 的重力沿斜面向下的分力为 $m_Ng\sin30^\circ=5\UN$，N 的最大静摩擦力等于滑动摩擦力 $f=\mu m_Ng\cos30^\circ=7.5\UN$。释放后 M 沿光滑斜面下滑并压缩弹簧，当弹簧弹力等于 N 的重力分力与最大静摩擦力之差 $2.5\UN$ 时，N 才开始运动。N 开始运动瞬间，弹簧弹力 $F=2.5\UN$，对 M 由牛顿第二定律 $m_Mg\sin30^\circ-F=m_Ma_M$，代入数据得 $a_M=0$，故 A 正确；

B．N 的速度最大时，其沿斜面方向合力为零，此时弹簧弹力仍为 $F=2.5\UN$，与 N 刚要运动瞬间弹簧弹力大小相等、形变量相同。N 开始运动后，对 M、N 和弹簧组成的系统，沿斜面方向的合力为 $F_{\text{合}}=m_Mg\sin\theta+m_Ng\sin\theta-f=0$，因此系统沿斜面方向动量守恒。从 N 开始运动到速度达到最大过程中，根据动量守恒有 $m_Mv_0=m_Mv_M+m_Nv_N$，根据能量守恒有 $(m_M+m_N)gx\sin\theta-\mu m_Ngx\cos\theta=\frac{1}{2}m_Mv_M^2+\frac{1}{2}m_Nv_N^2-\frac{1}{2}m_Mv_0^2$，解得 $v_M=-\frac{1}{3}v_0$（$v_0$ 为 M 在 N 开始运动瞬间的速度）。因此 N 的速度达到最大时，M 的速度不为零，故 B 错误；

C．N 开始运动后，在弹簧继续压缩的过程中，位移关系有 $x_M>x_N$，系统重力势能变化量与克服摩擦力做功的关系为 $m_Mgx_M\sin\theta+m_Ngx_N\sin\theta-\mu m_Ngx_N\cos\theta>0$，说明此过程系统重力势能变化量大于克服摩擦力做功，即重力对系统做功大于系统因摩擦而损失的能量，因此 M、N 的总动能与弹簧弹性势能之和增加，并不保持不变，故 C 错误；

D．第一次共速时弹簧处于压缩状态，第二次共速时弹簧处于伸长状态。两次共速之间 $x_N>x_M$，动量守恒使两次共速时的速度相同、系统动能相同。结合 C 项分析可知，$x_N>x_M$ 时系统重力势能变化量与克服摩擦力做功的关系为 $m_Mgx_M\sin\theta+m_Ngx_N\sin\theta-\mu m_Ngx_N\cos\theta<0$，说明此过程系统重力势能变化量小于克服摩擦力做功，即重力对系统做功小于系统因摩擦而损失的能量，因此 M、N 的总动能与弹簧弹性势能之和减少，而总动能不变，因此弹性势能减小，故 D 正确。

故选 AD。
}

\item
%% number: 8
%% paperName: 2026年福建高考第 8 题 6 分
%% typeId: 2
%% type: 多选题
%% chapter: 电磁感应
%% point: 电磁感应含容电路
%% method: 无
%% score: 6
%% degree: 251
%% duplicateId: 0
%% body:
如图，两条相距 $L$ 的光滑平行金属导轨位于同一水平面内，左端接一电容为 $C$ 的电容器。整个装置处于竖直向下的匀强磁场中，磁感应强度大小为 $B$。质量为 $m$ 的金属棒垂直导轨放置。初始时，棒以速度 $v_0$ 向右做匀速直线运动，到达 $MN$ 位置瞬间，开始对棒施加一与速度方向相反的外力，外力的大小与速度的大小成正比，比例系数为 $k$，棒最终停在 $PQ$ 位置。运动过程中，棒始终与导轨垂直且接触良好，回路电阻忽略不计。则\xzanswer{BC}
\onepicture[w=7.12cm]{figs/08.png}
\fourchoices[answer=BC]
{棒越过 $MN$ 瞬间的加速度大小为 $\frac{kv_0}{m}$}
{$MN$ 与 $PQ$ 的间距为 $\frac{(m+B^2L^2C)v_0}{k}$}
{棒运动到 $MN$ 与 $PQ$ 正中间位置时，电容器存储的电荷量为 $\frac{1}{2}BLCv_0$}
{若磁感应强度大小为 $2B$，棒不能到达 $PQ$ 位置}
%% answer: BC
%% memo:
\memoanswer{
A．取向右为正方向，对金属棒由牛顿第二定律得 $-kv-BIL=ma$。回路电流为电容器的充电电流，$I=\frac{\Delta Q}{\Delta t}$；回路电阻不计，电容器两端电压等于金属棒的动生电动势 $U=BLv$，故 $Q=CBLv$，$I=CBL\frac{\Delta v}{\Delta t}=CBLa$。代入得 $-(m+B^2L^2C)a=kv$，即 $a=-\frac{kv}{m+B^2L^2C}$。棒越过 $MN$ 瞬间速度为 $v_0$，加速度大小 $a=\frac{kv_0}{m+B^2L^2C}$，故 A 错误；

B．由 $a=\frac{\Delta v}{\Delta t}$ 及上式得 $-(m+B^2L^2C)\Delta v=kv\Delta t$，对全过程累加，初速度为 $v_0$、末速度为 $0$，故 $(m+B^2L^2C)v_0=kx$，即 $MN$ 与 $PQ$ 的间距 $x=\frac{(m+B^2L^2C)v_0}{k}$，故 B 正确；

C．由 $-(m+B^2L^2C)\Delta v=kv\Delta t$ 可得 $v=v_0-\frac{kx}{m+B^2L^2C}$，速度随位移均匀减小。运动到正中间位置时 $v=\frac{v_0}{2}$，电容器存储的电荷量 $Q=CBLv=\frac{1}{2}BLCv_0$，故 C 正确；

D．若磁感应强度大小为 $2B$，则总位移 $x'=\frac{(m+4B^2L^2C)v_0}{k}>x$，棒能越过原 $PQ$ 位置，故 D 错误。

故选 BC。
}

\item
%% number: 9
%% paperName: 2026年福建高考第 9 题 3 分
%% typeId: 3
%% type: 填空题
%% chapter: 光学
%% point: 光的折射
%% method: 无
%% score: 3
%% degree: 630
%% duplicateId: 0
%% body:
空气中，两块相同的平行玻璃板 Ⅰ 和 Ⅱ 镜像对称放置于水平底座上，玻璃板与底座的夹角均为 $\theta$（$\theta<90^\circ$），如图所示。一束绿光水平射入玻璃板 Ⅰ，入射光的延长线如图中虚线所示，则从玻璃板 Ⅰ 出射的光将出现在虚线\tkanswer[1.2cm]{上方}（填“上方”或“下方”），而从玻璃板 Ⅱ 出射的光与虚线重合。若将绿光换成红光，则从玻璃板 Ⅱ 出射的光\tkanswer[1.0cm]{会}（填“会”或“不会”）与虚线重合。（不考虑反射）
\onepicture[align=right, w=5.12cm]{figs/09.png}
%% answer:
\jdanswer{
上方，会
}
%% memo:
\memoanswer{
（1）光线斜射入平行玻璃板时，出射光与入射光方向平行，但会发生侧移。水平向右的绿光从空气斜射入玻璃板 Ⅰ 时，折射角小于入射角，光线向玻璃板的上侧偏折，因此从玻璃板 Ⅰ 出射的光出现在虚线上方。

（2）玻璃板 Ⅰ 和 Ⅱ 镜像对称，绿光经过 Ⅰ 的向上侧移会被玻璃板 Ⅱ 的向下侧移完全抵消，最终出射光与原虚线重合。换成红光时，根据对称性，玻璃板 Ⅰ 对红光的向上侧移量仍等于玻璃板 Ⅱ 对红光的向下侧移量，因此从玻璃板 Ⅱ 出射的光会与虚线重合。
}

\item
%% number: 10
%% paperName: 2026年福建高考第 10 题 3 分
%% typeId: 3
%% type: 填空题
%% chapter: 电场
%% point: 电场叠加
%% method: 对称法
%% score: 3
%% degree: 290
%% duplicateId: 0
%% body:
如图，长方形 $ABCD$ 的四个顶点固定有电荷量均为 $Q$ 的点电荷，其中 $A$、$C$ 点上的点电荷带负电，$B$、$D$ 点上的点电荷带正电。$AB$ 和 $BC$ 的长度分别为 $2a$ 和 $\sqrt3a$。$O$ 为长方形中心，长为 $2a$ 的线段 $OP$ 与长方形所在平面垂直，$M$ 为 $AB$ 边中点。将一电荷量为 $q$（$q>0$）的点电荷从 $P$ 点移到 $M$ 点，电场力对该点电荷所做的功\tkanswer[1.8cm]{等于零}（填“大于零”“等于零”或“小于零”），该点电荷在 $M$ 点所受电场力大小为\tkanswer[2.4cm]{$\frac{7kQq}{4a^2}$}。（静电力常量为 $k$）
\onepicture[align=right, w=5.20cm]{figs/10.png}
%% answer:
\jdanswer{
等于零，$\frac{7kQq}{4a^2}$
}
%% memo:
\memoanswer{
（1）平面 $OPM$ 是等量异种电荷 $A$、$B$ 连线的中垂面，也是 $C$、$D$ 连线的中垂面，该平面上各点电势均为零，即 $P$、$M$ 两点电势均为零。将一电荷量为 $q$ 的点电荷从 $P$ 点移到 $M$ 点，电场力对该点电荷所做的功等于零。

（2）$A$、$B$ 两点电荷在 $M$ 点的场强均为 $E_1=E_2=\frac{kQ}{a^2}$，方向沿 $BA$ 方向；$C$、$D$ 两点电荷到 $M$ 点的距离均为 $2a$，在 $M$ 点的场强均为 $E_3=E_4=\frac{kQ}{(2a)^2}=\frac{kQ}{4a^2}$，方向与 $AB$ 的夹角均为 $60^\circ$。$M$ 点的场强 $E=E_1+E_2-E_3\cos60^\circ-E_4\cos60^\circ=2\frac{kQ}{a^2}-2\cdot\frac{kQ}{4a^2}\cdot\frac{1}{2}=\frac{7kQ}{4a^2}$。该点电荷在 $M$ 点所受电场力大小为 $F=qE=\frac{7kQq}{4a^2}$。
}

\item
%% number: 11
%% paperName: 2026年福建高考第 11 题 3 分
%% typeId: 3
%% type: 填空题
%% chapter: 气体性质
%% point: 查理定律
%% method: 无
%% score: 3
%% degree: 401
%% duplicateId: 0
%% body:
如图，盛有液体的粗细均匀 U 形细管竖直放置，两端分别与装有一定质量理想气体的气瓶甲、乙相连，构成一个密闭系统。初始时，U 形管两侧液面高度相同，瓶内气体压强均为 $p_0$，温度均为 $T_0$。现分别对气瓶甲、乙缓慢加热，U 形管左侧液面下降，当两侧液柱高度差为 $h$（$h$ 较小，两侧气体体积变化可忽略）时，气瓶\tkanswer[0.8cm]{甲}（填“甲”或“乙”）内气体的温度较高，两气瓶内气体的温度差绝对值为\tkanswer[2.6cm]{$\frac{\rho ghT_0}{p_0}$}。已知重力加速度大小为 $g$，液体密度为 $\rho$ 且保持不变。
\onepicture[align=right, h=6.26cm]{figs/11.png}
%% answer:
\jdanswer{
甲，$\frac{\rho ghT_0}{p_0}$
}
%% memo:
\memoanswer{
（1）U 形管左侧液面下降，说明甲气瓶的压强较大。气瓶甲、乙内理想气体在升温过程中体积不变，根据理想气体状态方程 $\frac{pV}{T}=C$ 可知，气瓶甲内气体的温度较高。

（2）设加热后气瓶甲、乙内理想气体的压强分别为 $p_1$、$p_2$，温度分别为 $T_1$、$T_2$。根据查理定律，对甲有 $\frac{p_0}{T_0}=\frac{p_1}{T_1}$，对乙有 $\frac{p_0}{T_0}=\frac{p_2}{T_2}$。由题意可知 $p_1-p_2=\rho gh$。因此两气瓶内气体的温度差绝对值 $|T_1-T_2|=\frac{T_0}{p_0}(p_1-p_2)=\frac{\rho ghT_0}{p_0}$。
}

\item
%% number: 12
%% paperName: 2026年福建高考第 12 题 5 分
%% typeId: 4
%% type: 实验题
%% chapter: 运动和力
%% point: 动量定理
%% method: 无
%% score: 5
%% degree: 480
%% duplicateId: 0
%% body:
小组利用如图（a）所示的实验装置验证动量定理。所用器材有：气垫导轨，光电门，安装有遮光板的滑块，游标卡尺等。
\begin{enumerate}
	\item
	将气垫导轨的一端用垫块垫高。测得两光电门中心距离 $l$ 为 $50.00\Ucm$，两光电门处轨道的高度差 $h$ 为 $8.00\Ucm$，滑块（含遮光板）质量 $m$ 为 $200.0\Ug$；用游标卡尺测量遮光板宽度 $d$，示数如图（b）所示，读出 $d=$ \tkanswer[1.8cm]{10.0}$\Umm$。
	\item
	将滑块移至位置 $O$，由静止释放，测出遮光板通过光电门 1、2 的遮光时间以及通过两光电门的时间间隔，重复多次测量，算出它们的平均值 $\Delta t_1$、$\Delta t_2$、$t$ 分别为 $10.00\UmsT$、$6.25\UmsT$、$390.00\UmsT$。
	\item
	基于所测数据，并利用动量变化量大小的表达式 $\Delta p=$ \tkanswer[4.5cm]{$m\left(\frac{d}{\Delta t_2}-\frac{d}{\Delta t_1}\right)$}（用 $m$、$d$、$\Delta t_1$、$\Delta t_2$ 表示），可算出 $\Delta p=$ \tkanswer[1.8cm]{0.120}$\Ukgms$（结果保留 3 位有效数字）。
	\item
	重力加速度大小取 $9.80\Umsq$，根据测量数据，计算可得滑块从光电门 1 运动到光电门 2 过程中合外力冲量大小 $I=0.122\UNs$。由此可得，在误差允许范围内，物体所受合外力冲量大小等于动量变化量大小。
	\item
	关于本实验，下列说法正确的是\tkanswer[0.9cm]{B}。（单选，填正确答案标号）

	A．垫高导轨的目的是为了平衡滑块所受的阻力

	B．若选用宽度为 $\frac{d}{2}$ 的遮光板，瞬时速率测量值会更准确

	C．若滑块从不同位置静止释放，测得的遮光时间不变
\end{enumerate}
\onepicture[align=right, w=14.00cm]{\includesvg{figs/12}}
%% answer:
\jdanswer{
（1）$10.0\Umm$；

（3）$m\left(\frac{d}{\Delta t_2}-\frac{d}{\Delta t_1}\right)$，$0.120\Ukgms$；

（5）B。
}
%% memo:
\memoanswer{
（1）图中是 10 分度游标卡尺，精度为 $0.1\Umm$。主尺读数为 $10\Umm$，游标尺 0 刻度与主尺刻度对齐，游标尺读数为 $0\times0.1\Umm$，最终读数为 $10\Umm+0\times0.1\Umm=10.0\Umm$。

（3）用遮光板的平均速度近似代替滑块通过光电门的瞬时速度，滑块经过光电门 1、2 的速度分别为 $v_1=\frac{d}{\Delta t_1}$、$v_2=\frac{d}{\Delta t_2}$，因此动量变化量 $\Delta p=mv_2-mv_1=m\left(\frac{d}{\Delta t_2}-\frac{d}{\Delta t_1}\right)$。代入 $m=0.200\Ukg$、$d=0.0100\Um$、$\Delta t_1=10.00\UmsT$、$\Delta t_2=6.25\UmsT$，计算得 $v_1=1.00\Ums$、$v_2=1.60\Ums$，$\Delta p=0.120\Ukgms$。

（5）A．本实验中垫高导轨后，滑块的合外力是重力沿斜面向下的分力，是为了利用重力分力作为合外力，故 A 错误；

B．遮光板宽度越小，遮光时间越短，遮光过程的平均速度越接近滑块通过光电门的瞬时速度，瞬时速率的测量值更准确，故 B 正确；

C．滑块从不同位置静止释放，到达光电门的速度不同，遮光板宽度固定，遮光时间会随速度变化，不可能保持不变，故 C 错误。

故选 B。
}

\item
%% number: 13
%% paperName: 2026年福建高考第 13 题 0 分
%% typeId: 4
%% type: 实验题
%% chapter: 电路
%% point: 传感器
%% method: 无
%% score: 0
%% degree: 540
%% duplicateId: 0
%% body:
某小组利用电磁感应原理自制位移传感器。所用器材有：空芯线圈 $\mathrm{L}_1$（匝数 $N_1=1200$ 匝），空芯线圈 $\mathrm{L}_2$（匝数 $N_2=2300$ 匝，表面贴有刻度尺），学生电源，多用电表，开关 S，导线若干。
\begin{enumerate}
	\item
	$\mathrm{L}_1$、$\mathrm{L}_2$ 同轴嵌套，用学生电源给 $\mathrm{L}_1$ 提供交流电，用多用电表测量 $\mathrm{L}_2$ 上的感应电压。改变 $\mathrm{L}_2$ 位置，感应电压值会随之改变。根据上述原理，将图（a）中的实物电路连线补充完整\_\_\_\_\_\_\_\_。
	\drawpicanswer{\onepicture[align=right, w=10.87cm]{figs/13a.png}}{\onepicture[align=right, w=11.85cm]{figs/13_an1.png}}
	\item
	打开学生电源，将输出电压调至 $16\UV$ 挡；接通开关 S，用多用电表测得 $\mathrm{L}_1$ 两端电压 $U_1=16.1\UV$。
	\item
	使 $\mathrm{L}_2$ 刻度尺零刻度线与 $\mathrm{L}_1$ 的上平面对齐，用多用电表测量此初始位置 $\mathrm{L}_2$ 两端电压 $U_2$ 的值。测量时，遵循\tkanswer[1.8cm]{由大到小}（填“由大到小”或“由小到大”）的原则，进行量程选择。最终，选择开关和指针的位置如图（b）所示，读出 $U_2=$ \tkanswer[1.2cm]{3.9}$\UV$。经计算，$\frac{U_2}{U_1}$ 明显小于 $\frac{N_2}{N_1}$，主要原因是\tkanswer[1.2cm]{漏磁}（填“涡流损耗”或“漏磁”）。
	\onepicture[align=right, w=9.30cm]{figs/13b.png}
	\item
	将 $\mathrm{L}_2$ 沿轴向缓慢抽出，测量 7 组 $\mathrm{L}_2$ 位移 $x$ 及其对应的电压 $U_2$ 值。这 7 组数据点已在图（c）坐标纸上描出，请根据所描数据点绘制 $U_2-x$ 关系图线\_\_\_\_\_\_\_\_。
	\drawpicanswer{\onepicture[align=right, w=11.28cm]{figs/13c.png}}{\onepicture[align=right, w=13.20cm]{figs/13_an2.png}}
	\item
	由所绘的图线可知，在测量范围内，$\mathrm{L}_2$ 两端电压 $U_2$ 与其位移 $x$ 成\tkanswer[1.2cm]{线性}（填“线性”或“非线性”）关系。由此，利用该传感器可将位移测量转换为电压测量。
\end{enumerate}
%% answer:
\jdanswer{
（1）见题图；

（3）由大到小，$3.9\UV$，漏磁；

（4）见图线；

（5）线性。
}
%% memo:
\memoanswer{
（1）按图连接电路，使学生电源交流输出与线圈 $\mathrm{L}_1$ 构成回路，多用电表交流电压挡接线圈 $\mathrm{L}_2$ 两端。

（3）为保护电表，量程选择应遵循由大到小的原则。选择开关和指针位置见题图（b），为交流 $10\UV$ 挡，最小刻度为 $0.2\UV$，读出 $U_2=3.9\UV$。经计算，$\frac{U_2}{U_1}$ 明显小于 $\frac{N_2}{N_1}$，主要原因是漏磁。

（4）舍去误差较大的点，让多数点落在直线上、少数点分布在直线两侧，连线见题图。

（5）由所绘图线可知，在测量范围内，$\mathrm{L}_2$ 两端电压 $U_2$ 与其位移 $x$ 成线性关系。
}

\item
%% number: 14
%% paperName: 2026年福建高考第 14 题 11 分
%% typeId: 6
%% type: 计算题
%% chapter: 运动和力
%% point: 动量和能量
%% method: 无
%% score: 11
%% degree: 770
%% duplicateId: 0
%% body:
如图，运动员在水平冰道上练习冰壶投掷时，将冰壶 A 以初速度 $v_0$ 从 $M$ 点掷出，使其沿直线滑行一段距离后停在 $N$ 点。随后，运动员又将冰壶 B 从 $M$ 点掷出，使其与冰壶 A 发生正碰。碰撞后，冰壶 A 最终停在大本营中心 $O$ 点。已知 $M$、$N$、$O$ 点在同一直线上，$N$、$O$ 两点距离为 $L$，冰壶与冰道间的动摩擦因数为 $\mu$，重力加速度大小为 $g$；冰壶 A、B 质量相同，均视为质点，碰撞时间极短，可视为弹性碰撞。求：
\begin{enumerate}
	\item
	冰壶 A 从 $M$ 点滑行到 $N$ 点所用时间；
	\item
	碰撞后瞬间，冰壶 A 的速度大小；
	\item
	碰撞前瞬间，冰壶 B 的速度大小。
\end{enumerate}
\onepicture[align=right, w=12.97cm]{figs/14.png}
%% answer:
\jdanswer{
\begin{enumerate}
	\item
	$\frac{v_0}{\mu g}$

	\item
	$\sqrt{2\mu gL}$

	\item
	$\sqrt{2\mu gL}$
\end{enumerate}
}
%% memo:
\memoanswer{
（1）冰壶 A 在冰道上滑行时，滑动摩擦力提供加速度，由牛顿第二定律 $\mu mg=ma$，得加速度大小 $a=\mu g$。冰壶 A 做匀减速直线运动到 $N$ 点速度为零，由 $0=v_0-at$ 解得 $t=\frac{v_0}{\mu g}$。

（2）碰撞后冰壶 A 从碰撞位置滑行距离 $L$ 后停在 $O$ 点，滑行过程只有摩擦力做功，由动能定理 $-\mu mgL=0-\frac{1}{2}mv_A^2$，解得碰撞后 A 的速度 $v_A=\sqrt{2\mu gL}$。

（3）两冰壶质量相同，发生弹性正碰。碰撞前 A 静止，由动量守恒和能量守恒可知，碰撞后 B 静止、A 获得 B 碰前的速度，故碰撞前瞬间 B 的速度 $v_{B0}=v_A=\sqrt{2\mu gL}$。
}

\item
%% number: 15
%% paperName: 2026年福建高考第 15 题 12 分
%% typeId: 6
%% type: 计算题
%% chapter: 磁场
%% point: 带电粒子在组合场中的运动
%% method: 无
%% score: 12
%% degree: 310
%% duplicateId: 0
%% body:
如图（a），两块足够大的平行金属板水平放置，板间加有大小随时间 $t$ 变化的匀强电场和匀强磁场，变化规律分别如图（b）和（c）所示，电场强度幅值为 $E_0$，方向竖直向下，磁感应强度幅值为 $\frac{\sqrt3E_0}{3v_0}$，方向平行板面且垂直纸面向里。$t=0$ 时，一带正电的粒子以速率 $v_0$ 从 $M$ 点沿两金属板中轴线 $MN$ 向右射入；$t=t_0$ 时，粒子速度方向与 $MN$ 夹角为 $60^\circ$。运动过程中粒子始终不与板相碰，重力不计，忽略电场、磁场突变瞬间对粒子速度的影响。
\begin{enumerate}
	\item
	求粒子比荷；
	\item
	求 $t_0\sim3t_0$ 内粒子的速度大小和运动轨道半径；
	\item
	现调节 $0\sim8t_0$ 内磁感应强度幅值，使粒子在 $t=8t_0$ 时运动到中轴线，此时速度大小为 $v_0$，方向水平向右，且在 $t_0\sim3t_0$ 内粒子速度方向的偏转角小于 $360^\circ$，则幅值应调为多少？
\end{enumerate}
\onepicture[align=right, w=14.50cm]{figs/15.png}
%% answer:
\jdanswer{
\begin{enumerate}
	\item
	$\frac{\sqrt3v_0}{E_0t_0}$

	\item
	$2v_0$，$2v_0t_0$

	\item
	$\frac{\sqrt3\pi E_0}{9v_0}$ 或 $\frac{\sqrt3\pi E_0}{6v_0}$
\end{enumerate}
}
%% memo:
\memoanswer{
（1）$0\sim t_0$ 时间内粒子在电场力作用下运动。设粒子的电荷量为 $q$、质量为 $m$，加速度大小为 $a$，$t=t_0$ 时粒子速度在竖直方向的分量大小为 $v_y$。由牛顿运动定律和匀变速直线运动规律有 $qE_0=ma$，$v_y=at_0$。由速度的矢量分解关系得 $v_y=v_0\tan60^\circ=\sqrt3v_0$。联立得 $\frac{q}{m}=\frac{\sqrt3v_0}{E_0t_0}$。

（2）$t_0\sim3t_0$ 时间内粒子在磁场中做匀速圆周运动。粒子速度大小 $v_1=\frac{v_0}{\cos60^\circ}=2v_0$。设轨道半径为 $r$，磁感应强度幅值为 $B_0=\frac{\sqrt3E_0}{3v_0}$，由牛顿运动定律有 $qv_1B_0=\frac{mv_1^2}{r}$，代入数据得 $r=2v_0t_0$。

（3）设调整后磁感应强度幅值为 $B_1$，粒子在磁场中运动的周期 $T_1=\frac{2\pi m}{qB_1}$。粒子在 $t_0\sim3t_0$ 和 $5t_0\sim7t_0$ 内均在磁场中运动，每次运动时间均为 $2t_0$，速度方向偏转角 $\alpha=\frac{2t_0}{T_1}\cdot2\pi=2t_0\cdot\frac{qB_1}{m}$。

经分析，符合题意的两种情况为：

情况一：每次在磁场中速度方向偏转 $120^\circ$，则 $2t_0\cdot\frac{qB_1}{m}=\frac{2\pi}{3}$，解得 $B_1=\frac{\pi m}{3qt_0}=\frac{\sqrt3\pi E_0}{9v_0}$。

情况二：每次在磁场中速度方向偏转 $180^\circ$，则 $2t_0\cdot\frac{qB_1}{m}=\pi$，解得 $B_1=\frac{\pi m}{2qt_0}=\frac{\sqrt3\pi E_0}{6v_0}$。

故幅值应调为 $\frac{\sqrt3\pi E_0}{9v_0}$ 或 $\frac{\sqrt3\pi E_0}{6v_0}$。
}

\item
%% number: 16
%% paperName: 2026年福建高考第 16 题 16 分
%% typeId: 6
%% type: 计算题
%% chapter: 运动和力
%% point: 传送带
%% method: 无
%% score: 16
%% degree: 170
%% duplicateId: 0
%% body:
如图（a），物块 A 置于水平传送带上，通过跨过光滑轻质定滑轮的轻绳与右边的物块 B 相连。B 通过轻绳与物块 C 相连，连接各物块的轻绳均处于拉紧状态且不可伸长，A 与滑轮间的绳子与传送带平行。传送带始终以恒定速率顺时针转动。用挡板 P 使 A 保持静止。$t=0$ 时，撤去挡板 P，同时在 A 上施加一水平外力 $F$，A、B、C 从静止开始运动。以 $x$ 表示 A 离开初始位置的位移，$F$ 与 $x$ 成正比关系，取水平向右为 $F$ 和 $x$ 的正方向，变化关系如图（b）所示。A 第一次向右运动过程中，其速度 $v$ 随位移 $x$ 的变化关系如图（c）所示。当 $x=1\Um$ 时，A 与传送带共速。已知 A、B 质量分别为 $2\Ukg$、$3\Ukg$，A 与传送带间的动摩擦因数小于 $1$，最大静摩擦力等于滑动摩擦力，重力加速度大小取 $10\Umsq$，A 始终在传送带上，B、C 始终未与滑轮和地面相碰。
\begin{enumerate}
	\item
	求 A 离开初始位置向右运动 $1\Um$ 的过程中，$F$ 做的功；
	\item
	求 C 的质量和 A 与传送带间的动摩擦因数；
	\item
	当 A 第一次运动到 $x=6\Um$ 时，剪断 B、C 间的绳子，A 开始向左运动。求 A 第一次运动到最左端时相对 $x=0$ 位置的位移。
\end{enumerate}
\onepicture[align=right, w=14.50cm]{figs/16.png}
%% answer:
\jdanswer{
\begin{enumerate}
	\item
	$-10\UJ$

	\item
	$3\Ukg$，$0.5$

	\item
	$\frac{1}{2}-\frac{7\sqrt{10}}{10}\Um$
\end{enumerate}
}
%% memo:
\memoanswer{
（1）设 A、B、C 的质量分别为 $m_A$、$m_B$、$m_C$，A 与传送带间的动摩擦因数为 $\mu$，A 受到的滑动摩擦力大小为 $f$，A 从初始位置向右运动 $1\Um$、$4\Um$、$6\Um$ 的位移分别记为 $x_1$、$x_2$、$x_3$。依题意，A 与传送带间的动摩擦因数小于 $1$，经分析可知，A 向右运动过程中绳子处于拉紧状态。设 A 位移为 $x_1$ 时外力为 $F_1$，此过程中外力做的功 $W_1=\frac{0+F_1}{2}(x_1-0)$。由图（b）可知 $F_1=-20\UN$，代入得 $W_1=-10\UJ$。

（2）设 A 位移为 $x_3$ 时外力为 $F_3$，A 从初始位置运动到 $x=6\Um$ 过程中摩擦力做功为 $W_f$，则 $f=\mu m_Ag$，$W_f=f(x_1-0)-f(x_2-x_1)+f(x_3-x_2)$。对 A、B、C 组成的系统，从初始位置运动到 $x=6\Um$ 过程中，由动能定理有 $W_f+(m_B+m_C)g(x_3-0)+\frac{0+F_3}{2}(x_3-0)=0$。由图（b）可知 $F_3=-120\UN$，联立并代入数据得 $m_C=3\Ukg$。

设 A 位移为 $x_2$ 时外力为 $F_2$，对 A、B、C 组成的系统，从 $x=1\Um$ 运动到 $x=4\Um$ 过程中，由动能定理有 $-f(x_2-x_1)+(m_B+m_C)g(x_2-x_1)+\frac{F_1+F_2}{2}(x_2-x_1)=0$。由图（b）可知 $F_2=-80\UN$，联立并代入数据得 $\mu=0.5$。

（3）剪断 B、C 间的绳子后，A、B 一起运动。假设 A 向左运动过程中存在绳子拉力为零的位置，设绳子拉力刚为零时 A 的位移为 $x_4$，此时 A 的加速度大小为 $g$、方向向右，外力为 $F_4$，由牛顿第二定律有 $f+F_4=m_Ag$。由图（b）可知 $F_4=-20x_4$，联立并代入数据得 $x_4=-0.5\Um$。

设 A 位移为 $x_4$ 时 A、B 的速度大小为 $v$，A 从 $x=6\Um$ 运动到 $x=-0.5\Um$ 过程中，对 A、B 组成的系统，由动能定理有 $f(x_4-x_3)+m_Bg(x_4-x_3)+\frac{F_3+F_4}{2}(x_4-x_3)=\frac{1}{2}(m_A+m_B)v^2-0$。联立并代入数据得 $\frac{1}{2}m_Av^2=39\UJ$，可知 A 在 $x=-0.5\Um$ 时动能大于零，假设成立。

设 A 第一次运动到最左端的位移为 $x_5$，A 从 $x_4$ 处运动到 $x_5$ 处过程中，A、B 初速度相同，A 做加速度增大的减速运动，B 做加速度大小恒为 $g$ 的减速运动，故此过程 A、B 间绳子始终保持松弛。设 A 位移为 $x_5$ 时的外力为 $F_5$，由动能定理有 $f(x_5-x_4)+\frac{F_4+F_5}{2}(x_5-x_4)=0-\frac{1}{2}m_Av^2$。由图（b）可知 $F_5=-20x_5$，联立并代入数据得 $x_5=\frac{1}{2}\pm\frac{7\sqrt{10}}{10}\Um$。由于正根不符合实际，故 $x_5=\frac{1}{2}-\frac{7\sqrt{10}}{10}\Um$。
}

\end{enumerate}

\end{document}
